% GENERATED FILE. Edit canonical lessons, then run npm run generate:books.
% canonical-source-hash: fnv1a64:961ef255739a76fe
% canonical-lessons: SA-C178-matulani, SA-C178-svasurah, SA-C178-svasruh, SA-C178-jamata, SA-C178-snusa

\chapter{Mother's brother's wife, Father-in-law, Mother-in-law, Son-in-law, Daughter-in-law}
\label{ch:sa-mother-s-brother-s-wife-father-in-law-mother-in-law-son-in-law-daughter-in-law}
\hlchaptermodality{178}

\begin{chapteropening}
\textbf{By the end of this chapter:} \emph{I can say a mother's brother's wife, a father-in-law, a mother-in-law, a son-in-law and a daughter-in-law.}

Complete the last lesson of chapter 178: I can say a mother's brother's wife, a father-in-law, a mother-in-law, a son-in-law and a daughter-in-law.
\end{chapteropening}

\section[\texorpdfstring{\sk{मातुलानी} (mātulānī)}{मातुलानी (mātulānī)}]{\sk{मातुलानी} (mātulānī) — a mother's brother's wife}

\label{lesson:SA-C178-matulani}



\begin{quote}
Before the new one: say the Sanskrit for trusts, then the Sanskrit for is ashamed.
\end{quote}

\subsection*{You'll want to know: \sk{मातुलानी}}
\textbf{\sk{मातुलानी}} — \emph{mātulānī} — \textquotedblleft{}a mother's brother's wife\textquotedblright{}.

\begin{grammarlens}[title={What you've built}]
One more word for this chapter.
\end{grammarlens}

\subsection*{Your turn}
\begin{itemize}
\raggedright
  \item \emph{Say it:} \emph{mātulānī}
  \item \emph{Say it:} \emph{mātulānī}, once more
  \item \emph{Say it:} say \emph{mātulānī}
  \item \emph{Recall:} say the Sanskrit for trusts, then the Sanskrit for is ashamed, then say \emph{mātulānī} again
\end{itemize}

\begin{tcolorbox}[breakable,colback=teal!4,colframe=teal!35!black,title={Before you move on}]
What does \sk{मातुलानी} mean? (\textquotedblleft{}A mother's brother's wife\textquotedblright{}.) Say it once more.
\end{tcolorbox}
\section[\texorpdfstring{\sk{श्वशुरः} (śvaśuraḥ)}{श्वशुरः (śvaśuraḥ)}]{\sk{श्वशुरः} (śvaśuraḥ) — a father-in-law}

\label{lesson:SA-C178-svasurah}



\begin{quote}
Before the new one: say the Sanskrit for is ashamed, then the Sanskrit for a mother's brother's wife.
\end{quote}

\subsection*{You'll want to know: \sk{श्वशुरः}}
\textbf{\sk{श्वशुरः}} — \emph{śvaśuraḥ} — \textquotedblleft{}a father-in-law\textquotedblright{}.

\begin{grammarlens}[title={What you've built}]
One more word for this chapter.
\end{grammarlens}

\subsection*{Your turn}
\begin{itemize}
\raggedright
  \item \emph{Say it:} \emph{śvaśuraḥ}
  \item \emph{Say it:} \emph{śvaśuraḥ}, once more
  \item \emph{Say it:} say \emph{śvaśuraḥ}
  \item \emph{Recall:} say the Sanskrit for is ashamed, then the Sanskrit for a mother's brother's wife, then say \emph{śvaśuraḥ} again
\end{itemize}

\begin{tcolorbox}[breakable,colback=teal!4,colframe=teal!35!black,title={Before you move on}]
What does \sk{श्वशुरः} mean? (\textquotedblleft{}A father-in-law\textquotedblright{}.) Say it once more.
\end{tcolorbox}
\section[\texorpdfstring{\sk{श्वश्रूः} (śvaśrūḥ)}{श्वश्रूः (śvaśrūḥ)}]{\sk{श्वश्रूः} (śvaśrūḥ) — a mother-in-law}

\label{lesson:SA-C178-svasruh}



\begin{quote}
Before the new one: say the Sanskrit for a mother's brother's wife, then the Sanskrit for a father-in-law.
\end{quote}

\subsection*{You'll want to know: \sk{श्वश्रूः}}
\textbf{\sk{श्वश्रूः}} — \emph{śvaśrūḥ} — \textquotedblleft{}a mother-in-law\textquotedblright{}.

\begin{grammarlens}[title={What you've built}]
One more word for this chapter.
\end{grammarlens}

\subsection*{Your turn}
\begin{itemize}
\raggedright
  \item \emph{Say it:} \emph{śvaśrūḥ}
  \item \emph{Say it:} \emph{śvaśrūḥ}, once more
  \item \emph{Say it:} say \emph{śvaśrūḥ}
  \item \emph{Recall:} say the Sanskrit for a mother's brother's wife, then the Sanskrit for a father-in-law, then say \emph{śvaśrūḥ} again
\end{itemize}

\begin{tcolorbox}[breakable,colback=teal!4,colframe=teal!35!black,title={Before you move on}]
What does \sk{श्वश्रूः} mean? (\textquotedblleft{}A mother-in-law\textquotedblright{}.) Say it once more.
\end{tcolorbox}
\section[\texorpdfstring{\sk{जामाता} (jāmātā)}{जामाता (jāmātā)}]{\sk{जामाता} (jāmātā) — a son-in-law}

\label{lesson:SA-C178-jamata}



\begin{quote}
Before the new one: say the Sanskrit for a father-in-law, then the Sanskrit for a mother-in-law.
\end{quote}

\subsection*{You'll want to know: \sk{जामाता}}
\textbf{\sk{जामाता}} — \emph{jāmātā} — \textquotedblleft{}a son-in-law\textquotedblright{}.

\begin{grammarlens}[title={What you've built}]
One more word for this chapter.
\end{grammarlens}

\subsection*{Your turn}
\begin{itemize}
\raggedright
  \item \emph{Say it:} \emph{jāmātā}
  \item \emph{Say it:} \emph{jāmātā}, once more
  \item \emph{Say it:} say \emph{jāmātā}
  \item \emph{Recall:} say the Sanskrit for a father-in-law, then the Sanskrit for a mother-in-law, then say \emph{jāmātā} again
\end{itemize}

\begin{tcolorbox}[breakable,colback=teal!4,colframe=teal!35!black,title={Before you move on}]
What does \sk{जामाता} mean? (\textquotedblleft{}A son-in-law\textquotedblright{}.) Say it once more.
\end{tcolorbox}
\section[\texorpdfstring{\sk{स्नुषा} (snuṣā)}{स्नुषा (snuṣā)}]{\sk{स्नुषा} (snuṣā) — a daughter-in-law}

\label{lesson:SA-C178-snusa}



\begin{quote}
Before the new one: say the Sanskrit for a mother-in-law, then the Sanskrit for a son-in-law.
\end{quote}

\subsection*{You'll want to know: \sk{स्नुषा}}
\textbf{\sk{स्नुषा}} — \emph{snuṣā} — \textquotedblleft{}a daughter-in-law\textquotedblright{}.

\begin{grammarlens}[title={What you've built}]
One more word for this chapter.
\end{grammarlens}

\subsection*{Your turn}
\begin{itemize}
\raggedright
  \item \emph{Say it:} \emph{snuṣā}
  \item \emph{Say it:} \emph{snuṣā}, once more
  \item \emph{Say it:} say \emph{snuṣā}
  \item \emph{Recall:} say the Sanskrit for a mother-in-law, then the Sanskrit for a son-in-law, then say \emph{snuṣā} again
\end{itemize}

\begin{tcolorbox}[breakable,colback=teal!4,colframe=teal!35!black,title={Before you move on}]
What does \sk{स्नुषा} mean? (\textquotedblleft{}A daughter-in-law\textquotedblright{}.) Say it once more.
\end{tcolorbox}
